\section{研究動機}

FPGA 具有平行處理與即時控制的特性，適合用於音訊訊號處理與硬體控制。本專題希望利用 FPGA 製作一套可以錄製聲音、播放聲音，並即時顯示頻譜變化的系統。透過 Audio CODEC 完成聲音輸入與輸出，再搭配 SRAM、FLASH 作為暫存與永久儲存記憶體，最後利用 FFT 分析聲音頻率，使 LED 矩陣能隨著音樂頻率變化而跳動。藉由此專題，可以學習 FPGA 在音訊處理、記憶體控制、顯示介面與系統整合上的應用。

\section{研究目的}

\begin{enumerate}
    \item 使用 Audio CODEC WM8731 完成聲音輸入與輸出。
    \item 設定音訊取樣頻率為 22.4 kHz，資料格式為 16 bit。
    \item 錄音資料先暫存在 SRAM 中，確認後再寫入 FLASH。
    \item 播放時可從 SRAM 或 FLASH 讀取音訊資料並輸出聲音。
    \item 對播放中的聲音訊號進行 FFT 頻譜分析。
    \item 將 FFT 結果轉換成 8$\times$8 LED 矩陣顯示。
    \item 使用 BUTTON 或 SW 控制錄音、播放、儲存與停止功能。
    \item 使用 16$\times$2 LCD 顯示目前系統狀態與操作提示。
    \item 練習 FPGA 中 FSM、記憶體控制、音訊介面與顯示控制的整合設計。
\end{enumerate}

